% 配置性质、热反馈推导与数值约定；在本章小结之后、符号索引之前载入。
\clearpage
\section*{补充性质与证明}
\label{sec:ch3-properties}
\markright{补充性质与证明}

本部分仅适用于第3.6—3.7节的恒功率、二聚合热节点特例；不将配置能量分解或固定选项热行推广为第3.2—3.5节一般模型的性质。

\subsection*{配置能量与时序可行性}

本部分说明配置能量与时序可行性的关系，并给出热约束影响最优配置及能耗的条件。模式中的映射与频率决定执行时长、活动能量、资源冲突和注热位置；时序决定这些选择能否组成满足全部要求的完整安排。两类决策的联系如图\ref{fig:config-timing}所示。

\begin{figure}[!htbp]
\centering
\resizebox{\linewidth}{!}{%
\begin{tikzpicture}[node distance=7mm and 10mm,box/.style={draw,rounded corners,align=center,text width=4.5cm,minimum height=10mm,font=\small},>=Latex]
\node[box] (d) {Stage模式表与合法支持};
\node[box,below=of d] (m) {配置$m$：执行单元与频率};
\node[box,below left=of m] (e) {固定配置能量$E(m)$};
\node[box,below right=of m] (s) {时序$s$：执行顺序与等待};
\node[box,below=of s] (r) {资源、频率域与请求期限};
\node[box,below=of e] (p) {配置与时序共同确定\par 完整节点功率轨迹$p(t)$};
\node[box,below=of p] (t) {入口状态与RC传播\par 连续热安全};
\node[box,below=of r] (f) {是否存在完整可行安排};
\node[box,below right=of t] (o) {选择最低可实现能量};
\draw[->] (d)--(m); \draw[->] (m)--(e); \draw[->] (m)--(s);
\draw[->] (s)--(r); \draw[->] (m.west) -| ([xshift=-4mm]p.west) -- (p.west);
\draw[->] (s)--(p); \draw[->] (p)--(t); \draw[->] (r)--(f);
\draw[->] (t)--(f); \draw[->] (f)--(o);
\draw[->] (e.west) -- ++(-10mm,0) |- (o.west);
\end{tikzpicture}%
}

\caption{配置能量与时序可行性的关系。图示为问题结构，不规定一次性的先配置后调度算法。}
\label{fig:config-timing}
\end{figure}

在上述特例的固定窗口、恒定基准功率、固定模式响应和零额外开销的条件下，每个Stage在窗口内完整执行一次，等待与重排不改变配置的总能量。因此，可以将“配置的代价是多少”与“该配置是否存在安全安排”分别描述。对固定配置\(m\)，记满足依赖、资源和期限等无热约束的启动向量集合为\(\mathcal S_0(m)\)，进一步满足温度约束的集合为\(\mathcal S_T(m)\)。存在相应时序的配置分别组成
\[
\mathcal M_0=\{m:\mathcal S_0(m)\ne\varnothing\},\qquad
\mathcal M_T=\{m:\mathcal S_T(m)\ne\varnothing\},
\]
\noindent\begin{minipage}{\linewidth}
\setlength{\parindent}{2em}
其中，所有配置均取自\(\prod_iM_i\)。以下集合与最优值均在相同启动域下比较。对于连续版本，时序集合允许连续起点；对于第\ref{sec:milp}节的网格版本，则统一限制到给定启动网格。联合问题可等价写为

\begin{equation}
\boxed{\min_{\substack{m\in\prod_{i\in V}M_i\\s\in\mathcal S_T(m)}}E(m)
=\min_{m\in\mathcal M_T}E(m).}
\label{eq:ch3-16}
\end{equation}

\textbf{命题3-1：} 式\eqref{eq:ch3-16}与联合问题具有相同的可实现目标值集合。

\textbf{证明：} 任意联合可行解\((m,s)\)中的配置都属于\(\mathcal M_T\)，能量为\(E(m)\)；反之，若\(m\in\mathcal M_T\)，则至少存在一条完整安全时序\(s\)，使\((m,s)\)成为联合可行解，能量同样为\(E(m)\)。两种表达取得的目标值集合因此相同。模式集合有限，\(\mathcal M_T\)非空时可取得其中最低能量；为空时则无完整可行安排。证毕。
\end{minipage}\par\medskip

该性质说明，找到一条满足全部约束的完整安排，就能证明其配置能量可以实现，即正文中所称的安全见证。一条时序失败则只排除了这一安排，判断整个配置不可行仍需覆盖其全部可能时序或给出充分排除依据。

温度要求对最优能耗的影响也可由此说明。设无热问题可行，其最优值为\(E_0^*=\min_{m\in\mathcal M_0}E(m)\)，最优配置集为\(\mathcal M_0^*=\{m\in\mathcal M_0:E(m)=E_0^*\}\)；若热问题也可行，记最优值为\(E_T^*\)。由于\(\mathcal M_T\subseteq\mathcal M_0\)，有

\begin{equation}
E_T^*\ge E_0^*,\qquad
E_T^*=E_0^*\ \Longleftrightarrow\ \mathcal M_T\cap\mathcal M_0^*\ne\varnothing.
\label{eq:ch3-17}
\end{equation}

只要原最优配置中还保留一个热可行配置，就能通过它的安全时序达到原最优值；反之，达到同一能量的热安全方案也必然来自这一交集。因此，最优配置集合缩小并不必然提高最低能耗。只有全部原最优配置被排除、但仍存在其他热可行配置时，有限配置集合才保证最优能量严格提高；若所有配置均被排除，则问题不可行。


\subsection*{选项热响应的展开}
\label{supp:option-response}

正文式\eqref{eq:thermal-affine}将节点温度分为基准与选项活动温升，此处给出各项的计算表达。在固定入口状态、线性时不变RC和可加活动功率下，每个选项的热响应可以分别独立计算。记全部选项为\(\mathcal O=\bigcup_{i\in V}\mathcal O_i\)，用热传播核描述温度状态随时间的变化：
\begin{equation}
K(t)=\exp(-C_{\rm th}^{-1}Gt),\qquad t\ge0.
\label{eq:thermal-kernel}
\end{equation}
温度基准由入口状态和持续的空闲注热产生，选项温升则由对应执行区间内的活动注热产生。设\(p^0=[p_C^0,p_N^0]^{\mathsf T}\)，\(\theta_0=\theta(0)\)。CPU模式选项\(o=(i,m,k)\)的活动注热向量为\(v_o=[Q_i(m),0]^{\mathsf T}\)，NPU模式则为\([0,Q_i(m)]^{\mathsf T}\)，单位均为W。边界外公共功率只计入系统能量，不进入\(p^0\)或\(v_o\)。由式\eqref{eq:ch3-14}，基准温度与选项响应分别为
\begin{equation}
\begin{aligned}
b_n(t)&=T_a+\left[K(t)\theta_0+
\int_0^t K(t-u)C_{\rm th}^{-1}p^0\,\mathrm du\right]_n,\\
h_{no}(t)&=
\begin{cases}
\displaystyle\left[\int_{t_o}^{\min(t,e_o)}
K(t-u)C_{\rm th}^{-1}v_o\,\mathrm du\right]_n,&t>t_o,\\
0,&t\le t_o.
\end{cases}
\end{aligned}
\label{eq:thermal-option-response}
\end{equation}
其中，\(b_n\)为包含初温和固定空闲注热的绝对温度，单位为°C；\(h_{no}\)为活动温升，单位为K，其数值可与°C温差相加。选项尚未启动时，活动温升为零；执行结束后，已有温升仍随热传播变化。因此，\(h_{no}\)同时取决于注热节点、执行区间和观察时刻，与只表示功率积分的活动能量\(\varepsilon_o\)承担不同作用。

本特例的被动RC网络中，\(-C_{\rm th}^{-1}G\)的非对角元素非负，因而\(K(t)\)逐元素非负；结合\(v_o\ge0\)，可得\(h_{no}(t)\ge0\)。上述系数依赖选项响应的独立性；若加入温度相关泄漏，或并发改变选项时长与功率，则需重新审查这种表达。


\subsection*{热反馈的有效性与覆盖关系}
\label{supp:feedback-properties}

将热演化与调度联合写入MILP已有研究基础。Pierson等采用离散时间热模型建立资源分配与能量优化约束\cite{pierson2020}；本章则在连续耦合RC模型下，依据完整轨迹的违规位置生成约束。有限主问题与违规分离交替进行的组织方式可借鉴Subramanyam等的研究\cite{subramanyam2020}，约束有效性及终止条件仍需按本问题的假设推导，其收敛结论不能直接移用。

正文中的反馈行采用响应下界与放宽右端，并对源候选验证严格分离。下面分别说明这两项条件所保证的安全保留与排除范围。

\noindent\begin{minipage}{\linewidth}
\setlength{\parindent}{2em}
\textbf{命题3-2：} 式\eqref{eq:thermal-safe-row}保留全部原热安全安排。

\textbf{证明：} 若安排违反该行，由\(y_o\ge0\)及响应下界可得
\begin{equation}
T_n(\bar t;y)\ge\underline b_n+\sum_o a_oy_o
>\underline b_n+r\ge T_n^{\max}+\epsilon_T.
\label{eq:thermal-row-validity}
\end{equation}
故被排除安排必然超温，不属于原热安全集合。放宽右端只会减小排除范围，不会改变最终安全接纳标准。\hfill$\square$
\end{minipage}\par\medskip

\textbf{命题3-3：} 在满足逐Stage唯一选择的合法整数域中，若式\eqref{eq:thermal-source-separation}成立，则热行保留的安排集合包含于对应完整赋值排除保留的集合。

\textbf{证明：} 完整赋值排除只删除同时选择\(J^*\)中全部\(N\)个选项的安排，唯一选择约束保证该安排就是源安排。式\eqref{eq:thermal-source-separation}保证源安排也被热行删除，故满足热行的整数安排必满足对应完整赋值排除。\hfill$\square$

只有存在另一个同样违反该热行的合法整数安排时，上述包含关系才严格。命题说明违规温度信息可以用于排除其他安排，其范围限于满足唯一选择约束的合法整数域；分数解上的支配关系及实际求解加速仍需分别研究。

上述两项条件也给出有限域内的终止依据。假设主问题精确求解，连续热判定能够按原温度约束将每个候选明确分为安全或违规，不存在未分辨状态；若每轮新增约束全域有效，并严格排除当前违规候选，该候选便不会再次出现。选项域有限，因此在预算不限时，反馈经过有限轮次必得到安全安排或排空合法域。进一步，若每轮返回当前松弛问题的真正最优解，首次被接纳的安全解就是原有限网格热问题的最优解：松弛最优值提供下界，该安全解提供相同的可行上界。

实际求解的停止结论还取决于数值判定与计算预算。响应系数和连续峰温包络的数值保证均限于已核对的对称二节点实现；一般RC线性关系本身不提供任意网络的数值认证。求解过程分别保存安全\(U\)、独立\(L_0\)、求解器内部界及未决状态，有效热行也不独立认证求解器内部界。因此，超时、未分辨或主问题未收敛时，结果仍按已取得的证据报告，不能由有限域终止论证推定搜索已经完成。


\subsection*{数值复核与实现约定}
\label{supp:numerical-checks}
\enlargethispage{\baselineskip}

峰温误差包络与判定带宽承担不同作用：前者控制峰温计算误差，后者区分接近阈值的情况。对系数、右端与求和也需分别采用可靠包络，不能仅将温度容差视作计算误差保证。式\eqref{eq:thermal-safe-row}采用非负响应下界，放宽右端会减小排除范围，但不改变最终安全接纳标准。

完整轨迹已证违规后，名义峰时刻只用于定位。实现需在具体可表示时刻重新构造点证书，并对源候选的被选集合\(J^*\)计算可靠求和下界，验证式\eqref{eq:thermal-source-separation}。全局峰温下界不能代替某个时刻的点温度下界，普通浮点和略大于右端也不能作为严格分离证据。完整轨迹违规但点证书或分离条件未成立时，回退式\eqref{eq:no-good}；轨迹未分辨时，两种排除均不生成。

检查器先按指定整数容差解码近整数候选，再根据真实区间复核唯一选择、依赖、资源和频率域要求。整数结构约束与浮点热行分开检查，热行采用可靠区间求和。解码或结构复核失败记为数值异常，不接纳候选，也不据此生成热排除。

一次调用按完整选项集合识别缓存候选，可复用完全相同候选的热检查结果；加入新热行后，仍需重新检查其是否满足当前约束。热行按完整系数向量及右端的精确相等识别重复，只保留一条。重复行不构成新进展；再次返回已被热行明确排除的违规候选时，记录容差事件并改用该完整赋值的排除。若候选还违反已加入的整数排除约束，则停止为数值异常，不反复重启。

热行检查与完整轨迹检查保持一致：热行比较未分辨、完整轨迹已证安全时，可依据完整检查保存安全\(U\)；若达到\(U=L_0\)，仍可独立证明最优值。未达到独立上下界闭合且最终热行比较仍未分辨时，停止为数值未决。热行明确违规而完整轨迹被判安全时，记录矛盾并停止为异常。实现错误不作为正常回退处理。

有限网格的安全安排也是连续启动问题的安全安排；网格最优或不可行结论则限于该受限域。只有连续域有效下界被某个网格安全安排达到时，才能进一步确定连续域最优值。\(L_0\)的独立Stage构造对两种域均成立，使用时仍须确认各模式集非空，并区分最优值与全部最优配置集合。

